% =====================================================================
%  figset_A.tex  --  TikZ figure set A of the QCD factorization monograph
%  (agent figA, local style prefix "fa").
%  Scope (FIGURES_INDEX_v2): Sec. 2 -- Ch1 intro, Ch2 qcd, Ch3 prelim,
%  Ch5 sudakov, Ch6 ee;  Sec. 4.1 -- Ch2 qcd (required + optional).
%
%  Requirements: common/preamble_fact.tex (ctexbook, XeLaTeX).  It loads
%  tikz with the libraries decorations.pathmorphing, decorations.markings,
%  arrows.meta, calc, positioning, shapes.geometric, patterns, backgrounds,
%  fit, and pgfplots (compat=1.17), and contains the style blocks STYLES and
%  STYLES-EXT (quark, gluon, gluon loop, photon, ghost, mom, mom', vertex,
%  blob, wilson, wilson plain, cut, hard blob, jet blob, jetB blob, soft blob,
%  hadron blob, glauber, opvertex, scalar, longit + \longitmark, lcaxis,
%  region).  Every extra style is local to its tikzpicture and carries the
%  prefix "fa"; no package or library is loaded here.
%
%  Conventions used throughout (WRITER_SPEC_v2 Sec. 3):
%    * solid line with arrow = quark, the arrow is the fermion-number flow;
%    * dotted line with arrow = ghost (ghost-number flow);
%    * separate thin labelled arrows (mom/mom') = momentum flow;
%    * double line with arrow = Wilson line, arrow = path orientation
%      (= fermion-number flow of the parton it replaces);
%    * red dashed line = final-state cut;  zig-zag = branch cut;
%    * n_A, n_B = light-like directions of the two jets/beams (row 2).
%  fig:intro-dependency: node positions and edge routes were generated by
%  figures/test_fa/gen_dep.py from FINAL_OUTLINE_v2 (statement tags, the
%  \uses lists -> transitive reduction among the displayed results, and the
%  master-table column of bottleneck assumptions); the output is plain TikZ.
%  Its circled digits are typeset with a local \fontspec{FandolSong-Regular.otf}
%  (FandolSong ships with TeX Live; xeCJK treats U+2460-2464 as non-CJK, so
%  the Latin text font would not have the glyphs).  Rerun the script after a
%  change of labels or dependencies rather than editing coordinates by hand.
%  Each figure is delimited by  %%% BEGIN <label>  /  %%% END <label>.
% =====================================================================


%%% BEGIN fig:intro-dependency
\begin{figure}[htbp]
\centering
% Layout generated by agent figA from FINAL_OUTLINE_v2: node fill = rigour tag; grey arrows = transitive
% reduction of the \uses relation among the displayed results; circled digits = tools of column 1 used;
% orange dashed = numbered assumptions (premises of R_cond theorems); thick red = bottleneck assumptions.
\begin{tikzpicture}[
  fa res/.style={draw=#1!75!black,fill=#1!14,line width=0.6pt,rounded corners=2.5pt,
    text width=1.74cm,align=flush center,inner xsep=1pt,inner ysep=2pt,font=\fontsize{7.5}{8.8}\selectfont},
  fa cond/.style={dash pattern=on 2.2pt off 1.4pt},
  fa asm/.style={draw=orange!75!black,fill=orange!7,line width=0.6pt,rounded corners=2pt,
    dash pattern=on 2.2pt off 1.4pt,align=flush center,inner xsep=2.5pt,inner ysep=2pt,font=\scriptsize},
  fa dep/.style={draw=gray!55!black,line width=0.5pt,rounded corners=2.5pt,-{Stealth[length=3.4pt,width=3pt]}},
  fa bot/.style={draw=red!70!black,line width=1.2pt,rounded corners=3pt,-{Stealth[length=4.6pt,width=4.2pt]}},
  fa use/.style={draw=orange!75!black,line width=0.8pt,dash pattern=on 2.5pt off 1.5pt,rounded corners=3pt,
    -{Stealth[length=4pt,width=3.6pt]}},
  fa tool/.style={inner sep=0pt,font=\small},
  fa head/.style={align=center,font=\scriptsize,inner sep=1pt}]
  \node[fa head] at (0.92,1.42) {\textbf{工具}\\[1pt]{\tiny 第 2--3 章}};
  \node[fa head] at (3.08,1.42) {\textbf{红外结构}\\[1pt]{\tiny 第 4--5 章}};
  \node[fa head] at (5.47,1.42) {\textbf{严格核心}\\[1pt]{\tiny 第 6--8 章}};
  \node[fa head] at (7.86,1.42) {\textbf{因子化机器}\\[1pt]{\tiny 第 9 章与附录 N、O}};
  \node[fa head] at (10.32,1.42) {\textbf{过程}\\[1pt]{\tiny 第 10--15 章}};
  \node[fa head] at (12.71,1.42) {\textbf{边界}\\[1pt]{\tiny 第 18 章}};
  \node[fa head] at (15.05,1.42) {\textbf{前沿}\\[1pt]{\tiny 第 22--24 章}};
  \node[fa res=OliveGreen] (ST) at (0.92,0.00) {ST 恒等式\\[-0.4pt]\mbox{{\scriptsize{\fontspec{FandolSong-Regular.otf}①}}\,\tiny\color{gray!25!black}\cref{thm:prelim-slavnov-taylor}}};
  \node[fa res=OliveGreen] (BETA) at (0.92,-1.15) {$\beta$ 函数\\[-0.4pt]\mbox{{\scriptsize{\fontspec{FandolSong-Regular.otf}②}}\,\tiny\color{gray!25!black}\cref{prop:prelim-beta}}};
  \node[fa res=OliveGreen] (WIL) at (0.92,-2.42) {Wilson 线的\\动量空间规则\\[-0.4pt]\mbox{{\scriptsize{\fontspec{FandolSong-Regular.otf}③}}\,\tiny\color{gray!25!black}\cref{prop:prelim-wilson-rules}}};
  \node[fa res=OliveGreen] (CUT) at (0.92,-3.70) {切割规则\\[-0.4pt]\mbox{{\scriptsize{\fontspec{FandolSong-Regular.otf}④}}\,\tiny\color{gray!25!black}\cref{thm:prelim-cutting-rules}}};
  \node[fa res=OliveGreen] (COMP) at (0.92,-4.97) {复合算符\\重整化\\[-0.4pt]\mbox{{\scriptsize{\fontspec{FandolSong-Regular.otf}⑤}}\,\tiny\color{gray!25!black}\cref{thm:prelim-composite-renorm}}};
  \node[fa res=OliveGreen] (PC) at (3.08,0.00) {红外幂次计数\\[-0.4pt]\mbox{\tiny\color{gray!25!black}\cref{prop:ir-power-counting}}};
  \node[fa res=OliveGreen] (LAN) at (3.08,-1.15) {Landau 方程\\[-0.4pt]\mbox{\tiny\color{gray!25!black}\cref{thm:ir-landau-momentum}}};
  \node[fa res=RoyalBlue] (LR) at (3.08,-2.30) {领头区域\\[-0.4pt]\mbox{\tiny\color{gray!25!black}\cref{thm:ir-leading-regions}\,{\fontsize{5.5}{6}\selectfont{\fontspec{FandolSong-Regular.otf}①}}}};
  \node[fa res=RoyalBlue] (SUD) at (3.08,-3.56) {Sudakov\\形状因子\\[-0.4pt]\mbox{\tiny\color{gray!25!black}\cref{thm:sudakov-factorization}\,{\fontsize{5.5}{6}\selectfont{\fontspec{FandolSong-Regular.otf}③}}}};
  \node[fa res=OliveGreen] (PML) at (5.47,0.00) {PDF 的矩\\[-0.4pt]\mbox{\tiny\color{gray!25!black}\cref{thm:pdf-moments-local}\,{\fontsize{5.5}{6}\selectfont{\fontspec{FandolSong-Regular.otf}③⑤}}}};
  \node[fa res=OliveGreen] (DGL) at (5.47,-1.20) {DGLAP 方程\\[-0.4pt]\mbox{\tiny\color{gray!25!black}\cref{thm:pdf-dglap}\,{\fontsize{5.5}{6}\selectfont{\fontspec{FandolSong-Regular.otf}①②}}}};
  \node[fa res=OliveGreen] (EE) at (5.47,-2.30) {$R(s)$ 有限性\\[-0.4pt]\mbox{\tiny\color{gray!25!black}\cref{thm:ee-R-finite}\,{\fontsize{5.5}{6}\selectfont{\fontspec{FandolSong-Regular.otf}①②③}}}};
  \node[fa res=RoyalBlue] (OPE) at (5.47,-3.45) {欧氏 OPE\\[-0.4pt]\mbox{\tiny\color{gray!25!black}\cref{thm:dis-ope}\,{\fontsize{5.5}{6}\selectfont{\fontspec{FandolSong-Regular.otf}①⑤}}}};
  \node[fa res=OliveGreen] (MOM) at (5.47,-4.60) {DIS 矩因子化\\[-0.4pt]\mbox{\tiny\color{gray!25!black}\cref{thm:dis-moments}\,{\fontsize{5.5}{6}\selectfont{\fontspec{FandolSong-Regular.otf}③}}}};
  \node[fa res=OliveGreen,fa cond] (MNP) at (5.47,-5.86) {DIS 矩的\\非微扰形式\\[-0.4pt]\mbox{\tiny\color{gray!25!black}\cref{thm:dis-moments-np}\,{\fontsize{5.5}{6}\selectfont{\fontspec{FandolSong-Regular.otf}①②⑤}}}};
  \node[fa res=OliveGreen,fa cond] (SUBN) at (7.86,0.00) {条件减除定理\\[-0.4pt]\mbox{\tiny\color{gray!25!black}\cref{thm:appN-subtraction}}};
  \node[fa res=RoyalBlue] (SUB) at (7.86,-1.20) {减除定理\\[-0.4pt]\mbox{\tiny\color{gray!25!black}\cref{thm:method-subtraction}\,{\fontsize{5.5}{6}\selectfont{\fontspec{FandolSong-Regular.otf}⑤}}}};
  \node[fa res=RoyalBlue] (TPL) at (7.86,-2.30) {因子化模板\\[-0.4pt]\mbox{\tiny\color{gray!25!black}\cref{thm:method-template}}};
  \node[fa res=RoyalBlue] (GDEF) at (7.86,-3.45) {围道变形判据\\[-0.4pt]\mbox{\tiny\color{gray!25!black}\cref{prop:method-glauber-deformation}\,{\fontsize{5.5}{6}\selectfont{\fontspec{FandolSong-Regular.otf}④}}}};
  \node[fa res=OliveGreen,fa cond] (GLO) at (7.86,-4.60) {Glauber 相消\\[-0.4pt]\mbox{\tiny\color{gray!25!black}\cref{thm:appO-main}}};
  \node[fa res=RoyalBlue] (WARD) at (7.86,-5.75) {Ward 恒等式\\[-0.4pt]\mbox{\tiny\color{gray!25!black}\cref{lem:method-ward-multi}\,{\fontsize{5.5}{6}\selectfont{\fontspec{FandolSong-Regular.otf}①③}}}};
  \node[fa res=RoyalBlue] (DIS) at (10.32,-1.20) {DIS 因子化\\[-0.4pt]\mbox{\tiny\color{gray!25!black}\cref{thm:disall-factorization}}};
  \node[fa res=RoyalBlue] (SIA) at (10.32,-2.30) {$e^+e^-\to hX$\\[-0.4pt]\mbox{\tiny\color{gray!25!black}\cref{thm:coll-sia}}};
  \node[fa res=RoyalBlue] (DY) at (10.32,-3.45) {单举 DY\\[-0.4pt]\mbox{\tiny\color{gray!25!black}\cref{thm:dy-factorization}}};
  \node[fa res=RoyalBlue] (TMD) at (10.32,-4.60) {TMD 因子化\\[-0.4pt]\mbox{\tiny\color{gray!25!black}\cref{thm:tmd-factorization}}};
  \node[fa res=RoyalBlue] (THR) at (10.32,-5.75) {推力分布\\[-0.4pt]\mbox{\tiny\color{gray!25!black}\cref{thm:scet-thrust}}};
  \node[fa res=BrickRed] (BTMD) at (12.71,-2.58) {$hh\to HHX$\\TMD 因子化\\[-0.4pt]\mbox{\tiny\color{gray!25!black}\cref{thm:break-tmd}\,{\fontsize{5.5}{6}\selectfont{\fontspec{FandolSong-Regular.otf}③}}}};
  \node[fa res=Plum] (COL) at (12.71,-3.83) {有色末态\\共线因子化\\[-0.4pt]\mbox{\tiny\color{gray!25!black}\cref{conj:dy-coloured-final}}};
  \node[fa res=BrickRed] (BSP) at (12.71,-5.12) {类空严格\\共线因子化\\[-0.4pt]\mbox{\tiny\color{gray!25!black}\cref{prop:break-spacelike}}};
  \node[fa res=RoyalBlue] (PIN) at (15.05,0.00) {物理区\\钳夹判据\\[-0.4pt]\mbox{\tiny\color{gray!25!black}\cref{thm:front-collins-pinch}}};
  \node[fa res=RoyalBlue] (QTMD) at (15.05,-1.36) {准 TMD\\因子化\\[-0.4pt]\mbox{\tiny\color{gray!25!black}\cref{thm:lamet-quasi-tmd}\,{\fontsize{5.5}{6}\selectfont{\fontspec{FandolSong-Regular.otf}③⑤}}}};
  \node[fa res=RoyalBlue] (LAM) at (15.05,-2.58) {LaMET 匹配\\[-0.4pt]\mbox{\tiny\color{gray!25!black}\cref{thm:lattice-lamet}}};
  \node[fa res=RoyalBlue] (TO) at (15.05,-3.83) {软因子\\时序无关性\\[-0.4pt]\mbox{\tiny\color{gray!25!black}\cref{thm:yzliu-time-ordering}}};
  \begin{scope}[on background layer]
    \draw[fa dep] (WIL.south) -- (CUT.north);
    \draw[fa dep] (LAN.south) -- (LR.north);
    \draw[fa dep] (LR.south) -- (SUD.north);
    \draw[fa dep] (PML.south) -- (DGL.north);
    \draw[fa dep] (OPE.south) -- (MOM.north);
    \draw[fa dep] (SUBN.south) -- (SUB.north);
    \draw[fa dep] (SUB.south) -- (TPL.north);
    \draw[fa dep] (DY.south) -- (TMD.north);
    \draw[fa dep] (MNP.north) -- (MOM.south);
    \draw[fa dep] (GDEF.north) -- (TPL.south);
    \draw[fa dep] (GLO.north) -- (GDEF.south);
    \draw[fa dep] (WARD.north) -- (GLO.south);
    \draw[fa dep] (PC.west) -| (2.000,-1.000) |- (LR.west);
    \draw[fa dep] (TPL.east) -| (8.900,-1.700) |- (DIS.west);
    \draw[fa dep] (TPL.east) -- (SIA.west);
    \draw[fa dep] (TPL.east) -| (8.900,-2.900) |- ($(DY.west)+(0,0.11)$);
    \draw[fa dep] ($(GDEF.east)+(0,-0.11)$) -| (9.060,-4.700) |- (THR.west);
    \draw[fa dep] ($(DY.east)+(0,0.11)$) -| (11.500,-3.600) |- (COL.west);
    \draw[fa dep] (DGL.east) -| (6.500,-1.750) |- ($(TPL.west)+(0,0.11)$);
    \draw[white,line width=2pt] (4.530,-0.530) -- (6.590,-0.530);
    \draw[white,line width=2pt] (6.830,-0.530) -- (13.840,-0.530);
    \draw[white,line width=2pt] (6.620,-0.720) -- (13.650,-0.720);
    \draw[white,line width=2pt] (4.530,-1.760) -- (12.560,-1.760);
    \draw[white,line width=2pt] (4.530,-3.940) -- (6.680,-3.940);
    \draw[white,line width=2pt] (4.260,-4.110) -- (11.520,-4.110);
    \draw[fa dep] (LR.east) -| (4.410,-1.000) |- (6.710,-0.530) |- (SUBN.west);
    \draw[fa dep] (LR.east) -| (4.410,-1.000) |- (13.960,-0.530) |- (QTMD.west);
    \draw[fa dep] (LR.east) -| (4.410,-1.500) |- (12.710,-1.760) -| (BTMD.north);
    \draw[fa dep] (LR.east) -| (4.410,-3.400) |- (6.800,-3.940) |- ($(GLO.west)+(0,0.11)$);
    \draw[fa dep] (DGL.east) -| (6.500,-0.900) |- (13.770,-0.720) |- (LAM.west);
    \draw[fa dep] (SUD.east) -| (4.140,-3.900) |- (11.640,-4.110) |- (BSP.west);
    \draw[fa dep] (LAN.east) -| (4.270,0.820) -| (PIN.north);
  \end{scope}
  \node[fa asm] (A1) at (3.05,-7.42) {非微扰 OPE [C]\\[-0.4pt]{\tiny\color{gray!25!black}\cref{asm:intro-np-ope}}};
  \node[fa asm] (A2) at (6.45,-7.42) {Glauber 区 [A]\\[-0.4pt]{\tiny\color{gray!25!black}\cref{asm:appO-mixing,asm:appO-deformation,asm:appO-rewritten-cs}}};
  \node[fa asm] (A3) at (10.32,-7.42) {桥接假设 [C]\\[-0.4pt]{\tiny\color{gray!25!black}\cref{asm:intro-bridge}}};
  \node[fa asm] (A4) at (12.85,-7.42) {Collins 条件 [A]\\[-0.4pt]{\tiny\color{gray!25!black}\cref{asm:appN-collins}}};
  \begin{scope}[on background layer]
    \draw[fa use] (A1.north) -- (MNP.south);
    \draw[fa use] (A2.north -| 6.61,0) |- ($(GLO.west)+(0,-0.11)$);
    \draw[fa bot] (A2.east) -| (9.230,-5.000) |- ($(DY.west)+(0,-0.11)$);
    \draw[fa bot] (A2.east) -| (9.230,-5.000) |- (TMD.west);
    \draw[fa bot] (A3.north) -- (THR.south);
    \draw[fa bot] (A4.west) -| (11.380,-6.000) |- (DIS.east);
    \draw[fa bot] (A4.west) -| (11.380,-6.000) |- (SIA.east);
    \draw[fa bot] (A4.west) -| (11.380,-6.000) |- ($(DY.east)+(0,-0.11)$);
    \draw[fa bot] (A4.west) -| (11.380,-6.000) |- (TMD.east);
  \end{scope}
  \begin{scope}[shift={(14.15,-4.62)},font=\fontsize{7.5}{9}\selectfont]
    \filldraw[fill=OliveGreen!14,draw=OliveGreen!75!black,line width=0.5pt,rounded corners=1pt] (0,-0.095) rectangle ++(0.36,0.19);
    \node[anchor=west,inner sep=1pt] at (0.42,0.000) {R 严格};
    \filldraw[fill=RoyalBlue!14,draw=RoyalBlue!75!black,line width=0.5pt,rounded corners=1pt] (0,-0.445) rectangle ++(0.36,0.19);
    \node[anchor=west,inner sep=1pt] at (0.42,-0.350) {A 全阶论证};
    \filldraw[fill=BurntOrange!14,draw=BurntOrange!75!black,line width=0.5pt,rounded corners=1pt] (0,-0.795) rectangle ++(0.36,0.19);
    \node[anchor=west,inner sep=1pt] at (0.42,-0.700) {P 部分};
    \filldraw[fill=Plum!14,draw=Plum!75!black,line width=0.5pt,rounded corners=1pt] (0,-1.145) rectangle ++(0.36,0.19);
    \node[anchor=west,inner sep=1pt] at (0.42,-1.050) {C 猜想};
    \filldraw[fill=BrickRed!14,draw=BrickRed!75!black,line width=0.5pt,rounded corners=1pt] (0,-1.495) rectangle ++(0.36,0.19);
    \node[anchor=west,inner sep=1pt] at (0.42,-1.400) {V 破坏};
    \draw[gray!60!black,line width=0.5pt,dash pattern=on 2.2pt off 1.4pt,rounded corners=1pt] (0,-1.845) rectangle ++(0.36,0.19);
    \node[anchor=west,inner sep=1pt] at (0.42,-1.75) {条件定理};
    \node[anchor=west,inner sep=1pt] at (-0.03,-2.1) {{\fontspec{FandolSong-Regular.otf}①⑤}\ 所用工具};
    \draw[fa bot] (0,-2.45) -- (0.36,-2.45);
    \node[anchor=west,inner sep=1pt] at (0.42,-2.45) {瓶颈假设};
    \draw[fa use] (0,-2.8) -- (0.36,-2.8);
    \node[anchor=west,inner sep=1pt] at (0.42,-2.8) {前提};
  \end{scope}
\end{tikzpicture}
\caption{本书主要结论的逻辑依赖。框的颜色表示严格性标签（绿 R、蓝 A、橙 P、紫 C、红 V），虚线框为 $R_{\rm cond}$ 条件定理；灰色箭头是各条目依赖记录在所示结论之间的传递约化，引用编号后的圈码表示用到左列的哪几项工具（ST 为 Slavnov--Taylor 恒等式）。下方橙色虚线框为编号假设，其中“Glauber 区”框代表 Glauber--软混合、围道形变与改写形式的 Collins 条件三条假设；橙色虚线箭头指向以之为前提的条件定理，粗红线标出决定各全阶论证地位的瓶颈假设。完整的依赖图见附录~P。}
\label{fig:intro-dependency}
\end{figure}
%%% END fig:intro-dependency


%%% BEGIN fig:prelim-feynman-rules
\begin{figure}[htbp]
\centering
\begin{tikzpicture}[fd,
  fa idx/.style={font=\footnotesize,inner sep=2pt},
  fa rule/.style={font=\small},
  fa cap/.style={font=\footnotesize}]
  % ===================== row 1: propagators =====================
  % (a) quark
  \begin{scope}
    \draw[quark] (0,0) node[fa idx,left]{$j$} -- (2.4,0) node[fa idx,right]{$i$};
    \path[mom={$p$}] (0,0) -- (2.4,0);
    \node[fa rule] at (1.2,-0.95) {$\dfrac{\ii(\sla p+m)}{p^2-m^2+\ii0}\,\delta_{ij}$};
    \node[fa cap] at (1.2,-1.85) {(a)\ 夸克传播子};
  \end{scope}
  % (b) gluon, Feynman gauge
  \begin{scope}[xshift=5.3cm]
    \draw[gluon] (0,0) node[fa idx,left]{$\mu,a$} -- (2.4,0) node[fa idx,right]{$\nu,b$};
    \path[mom={$k$},mom sep=0.32cm] (0,0) -- (2.4,0);
    \node[fa rule] at (1.2,-0.95) {$\dfrac{-\ii g_{\mu\nu}\,\delta^{ab}}{k^2+\ii0}$};
    \node[fa cap] at (1.2,-1.85) {(b)\ 胶子传播子（Feynman 规范 $\xi_g=1$）};
  \end{scope}
  % (c) ghost
  \begin{scope}[xshift=10.6cm]
    \draw[ghost] (0,0) node[fa idx,left]{$a$} -- (2.4,0) node[fa idx,right]{$b$};
    \path[mom={$p$}] (0,0) -- (2.4,0);
    \node[fa rule] at (1.2,-0.95) {$\dfrac{\ii\,\delta^{ab}}{p^2+\ii0}$};
    \node[fa cap] at (1.2,-1.85) {(c)\ 鬼场传播子};
  \end{scope}
  % ===================== row 2: three-point vertices =====================
  % (d) quark-gluon vertex
  \begin{scope}[yshift=-3.55cm]
    \coordinate (v) at (1.45,0);
    \draw[quark] (0.35,-0.85) node[fa idx,left]{$j$} -- (v);
    \draw[quark] (v) -- (0.35,0.85) node[fa idx,left]{$i$};
    \draw[gluon] (v) -- (2.75,0) node[fa idx,right]{$\mu,a$};
    \node[vertex] at (v) {};
    \node[fa rule] at (1.2,-1.75) {$\ii g\,\gamma^\mu\,(t^a)_{ij}$};
    \node[fa cap] at (1.2,-2.62) {(d)\ 夸克—胶子顶点};
  \end{scope}
  % (e) ghost-gluon vertex
  \begin{scope}[xshift=5.3cm,yshift=-3.55cm]
    \coordinate (v) at (1.2,-0.3);
    \draw[ghost] (0,-0.3) node[fa idx,left]{$c$} -- (v);
    \draw[ghost] (v) -- (2.4,-0.3) node[fa idx,right]{$a$};
    \draw[gluon] (v) -- (1.2,0.9) node[fa idx,above]{$\mu,b$};
    \path[mom'={$p$}] (v) -- (2.4,-0.3);
    \node[vertex] at (v) {};
    \node[fa rule] at (1.2,-1.75) {$-g\,f^{abc}\,p^\mu$};
    \node[fa cap] at (1.2,-2.62) {(e)\ 鬼场—胶子顶点};
  \end{scope}
  % (f) three-gluon vertex, all momenta incoming
  \begin{scope}[xshift=11.8cm,yshift=-3.55cm]
    \coordinate (v) at (0,0);
    \draw[gluon] (v) -- (180:1.2) coordinate (ea) node[fa idx,left]{$\mu,a$};
    \draw[gluon] (v) -- (-56:1.0) coordinate (eb) node[fa idx,anchor=west]{$\nu,b$};
    \draw[gluon] (v) -- (56:1.0) coordinate (ec) node[fa idx,anchor=west]{$\rho,c$};
    \path[mom={$k$},mom sep=0.3cm] (ea) -- (v);
    \path[mom={$p$},mom sep=0.3cm] (eb) -- (v);
    \path[mom={$q$},mom sep=0.3cm] (ec) -- (v);
    \node[vertex] at (v) {};
    \node[fa rule] at (0,-1.82)
      {$\begin{aligned}gf^{abc}\big[&g^{\mu\nu}(k-p)^\rho+g^{\nu\rho}(p-q)^\mu\\[-1pt]&+g^{\rho\mu}(q-k)^\nu\big]\end{aligned}$};
    \node[fa cap] at (0,-2.62) {(f)\ 三胶子顶点（动量 $k,p,q$ 均流入）};
  \end{scope}
  % ===================== row 3: four-gluon vertex =====================
  \begin{scope}[xshift=2.3cm,yshift=-7.75cm]
    \coordinate (v) at (1.2,0);
    % short straight stubs keep the four coils from piling up at the vertex
    \foreach \a/\l/\an in {135/{\mu,a}/-45,225/{\nu,b}/45,315/{\rho,c}/135,45/{\sigma,d}/-135} {
      \draw[plain] (v) -- ++(\a:0.16);
      \draw[gluon] ($(v)+(\a:0.16)$) -- ++(\a:1.04) node[fa idx,anchor=\an]{$\l$};
    }
    \node[vertex] at (v) {};
    \node[fa rule,anchor=west] at (3.3,0)
      {$\begin{aligned}
        -\ii g^2\big[&f^{abe}f^{cde}\,(g^{\mu\rho}g^{\nu\sigma}-g^{\mu\sigma}g^{\nu\rho})\\
        &+f^{ace}f^{bde}\,(g^{\mu\nu}g^{\rho\sigma}-g^{\mu\sigma}g^{\nu\rho})\\
        &+f^{ade}f^{bce}\,(g^{\mu\nu}g^{\rho\sigma}-g^{\mu\rho}g^{\nu\sigma})\big]
      \end{aligned}$};
    \node[fa cap,anchor=west] at (3.3,-1.45) {(g)\ 四胶子顶点};
  \end{scope}
\end{tikzpicture}
\caption{本书使用的 Feynman 规则（协变规范，$D_\mu=\partial_\mu-\ii gA^a_\mu t^a$，与 Peskin--Schroeder 第 16 章一致；图中胶子传播子取 Feynman 规范 $\xi_g=1$，一般 $\xi_g$ 的形式见正文）。
夸克与鬼场传播子中的动量沿线上的箭头；夸克—胶子顶点 $(t^a)_{ij}$ 中的 $i$ 属于沿箭头离开顶点的夸克线；三胶子顶点的动量 $k,p,q$ 均流入顶点；
鬼场—胶子顶点中 $a$ 为出射鬼场（动量 $p$ 流出顶点），$c$ 为入射鬼场，$b$ 为胶子的色指标。每个闭合夸克圈与闭合鬼场圈各带因子 $-1$。}
\label{fig:prelim-feynman-rules}
\end{figure}
%%% END fig:prelim-feynman-rules


%%% BEGIN fig:qcd-vacuum-polarization
\begin{figure}[htbp]
\centering
\begin{tikzpicture}[fd,
  fa idx/.style={font=\footnotesize,inner sep=1.5pt},
  fa note/.style={font=\footnotesize,align=center},
  fa cap/.style={font=\footnotesize}]
  % common: external gluons (k, mu, a) -> loop -> (k, nu, b);  loop radius 0.6
  % ---------------- (a) quark loop ----------------
  \begin{scope}
    \draw[gluon] (-1.6,0) -- (-0.6,0);
    \draw[gluon] (0.6,0) -- (1.6,0);
    \draw[quark] (-0.6,0) arc[start angle=180,end angle=0,radius=0.6];
    \draw[quark] (0.6,0) arc[start angle=0,end angle=-180,radius=0.6];
    \foreach \x in {-0.6,0.6} \node[vertex] at (\x,0) {};
    \path[mom={$k$},mom sep=0.3cm] (-1.6,0) -- (-0.65,0);
    \path[mom={$k$},mom sep=0.3cm] (0.65,0) -- (1.6,0);
    \node[fa idx,anchor=north] at (-1.55,-0.2) {$\mu,a$};
    \node[fa idx,anchor=north] at (1.55,-0.2) {$\nu,b$};
    \node[fa note] at (0,0.98) {$\times\,n_f$};
    \node[fa cap] at (0,-1.35) {(a)\ 夸克圈};
  \end{scope}
  % ---------------- (b) gluon loop ----------------
  \begin{scope}[xshift=4.0cm]
    \draw[gluon] (-1.6,0) -- (-0.6,0);
    \draw[gluon] (0.6,0) -- (1.6,0);
    \draw[gluon] (-0.6,0) arc[start angle=180,end angle=0,radius=0.6];
    \draw[gluon] (0.6,0) arc[start angle=0,end angle=-180,radius=0.6];
    \foreach \x in {-0.6,0.6} \node[vertex] at (\x,0) {};
    \path[mom={$k$},mom sep=0.3cm] (-1.6,0) -- (-0.68,0);
    \path[mom={$k$},mom sep=0.3cm] (0.68,0) -- (1.6,0);
    \node[fa idx,anchor=north] at (-1.55,-0.2) {$\mu,a$};
    \node[fa idx,anchor=north] at (1.55,-0.2) {$\nu,b$};
    \node[fa note] at (0,1.05) {对称因子 $\tfrac12$};
    \node[fa cap] at (0,-1.35) {(b)\ 胶子圈};
  \end{scope}
  % ---------------- (c) ghost loop ----------------
  \begin{scope}[xshift=8.0cm]
    \draw[gluon] (-1.6,0) -- (-0.6,0);
    \draw[gluon] (0.6,0) -- (1.6,0);
    \draw[ghost] (-0.6,0) arc[start angle=180,end angle=0,radius=0.6];
    \draw[ghost] (0.6,0) arc[start angle=0,end angle=-180,radius=0.6];
    \foreach \x in {-0.6,0.6} \node[vertex] at (\x,0) {};
    \path[mom={$k$},mom sep=0.3cm] (-1.6,0) -- (-0.65,0);
    \path[mom={$k$},mom sep=0.3cm] (0.65,0) -- (1.6,0);
    \node[fa idx,anchor=north] at (-1.55,-0.2) {$\mu,a$};
    \node[fa idx,anchor=north] at (1.55,-0.2) {$\nu,b$};
    \node[fa note] at (0,0.98) {闭合鬼场圈 $\times(-1)$};
    \node[fa cap] at (0,-1.35) {(c)\ 鬼场圈};
  \end{scope}
  % ---------------- (d) tadpole ----------------
  \begin{scope}[xshift=12.0cm]
    % external legs and the loop leave the vertex with short straight stubs
    \draw[gluon] (-1.6,0) -- (-0.16,0);
    \draw[gluon] (0.16,0) -- (1.6,0);
    \draw[plain] (-0.16,0) -- (0.16,0);
    \draw[plain] (0,0) -- (128:0.17);
    \draw[plain] (0,0) -- (52:0.17);
    \draw[gluon] (128:0.17) to[out=128,in=180,looseness=1.25] (0,1.22)
                            to[out=0,in=52,looseness=1.25] (52:0.17);
    \node[vertex] at (0,0) {};
    \path[mom={$k$},mom sep=0.3cm] (-1.6,0) -- (-0.5,0);
    \path[mom={$k$},mom sep=0.3cm] (0.5,0) -- (1.6,0);
    \node[fa idx,anchor=north] at (-1.55,-0.2) {$\mu,a$};
    \node[fa idx,anchor=north] at (1.55,-0.2) {$\nu,b$};
    \node[fa note] at (0,1.68) {维数正规化中为零};
    \node[fa cap] at (0,-1.35) {(d)\ 四胶子顶点的蝌蚪图};
  \end{scope}
  % ---------------- the result ----------------
  \node[font=\small,anchor=north] at (6.0,-1.85)
    {$\delta Z_3=\dfrac{\alpha_s}{4\pi\varepsilon}\Big[\dfrac53\,C_A-\dfrac43\,T_Fn_f\Big]$
     \quad（Feynman 规范；一般规范中 $C_A$ 的系数为 $\tfrac{13}{6}-\tfrac{\xi_g}{2}$）};
  \node[font=\small,anchor=north] at (6.0,-3.05)
    {$\Longrightarrow$\ 与 $\delta Z_2$、$\delta Z_{1F}$ 合成：\
     $\delta Z_g=\delta Z_{1F}-\delta Z_2-\tfrac12\,\delta Z_3=-\dfrac{\alpha_s}{4\pi\varepsilon}\,\dfrac{\beta_0}{2}$，\
     $\beta_0=\tfrac{11}{3}C_A-\tfrac43T_Fn_f$};
\end{tikzpicture}
\caption{胶子传播子的单圈修正（Feynman 规范，外线动量 $k$ 自左向右流动）：(a) 夸克圈（对 $n_f$ 味求和），(b) 胶子圈，(c) 鬼场圈，
(d) 四胶子顶点的蝌蚪图（无标度积分，维数正规化中为零）。(b)(c) 单独都不满足横向性，三者之和正比于 $k^2g^{\mu\nu}-k^\mu k^\nu$；
紫外极点给出 $\delta Z_3$，与夸克场和顶点的重整化常数合成规范无关的 $\delta Z_g$ 与 $\beta_0=\frac{11}3C_A-\frac43T_Fn_f$（VF-PDF-14，VF-QCD-4）。}
\label{fig:qcd-vacuum-polarization}
\end{figure}
%%% END fig:qcd-vacuum-polarization


%%% BEGIN fig:qcd-running-coupling
\begin{figure}[htbp]
\centering
\begin{tikzpicture}
\begin{axis}[name=A,width=0.47\textwidth,height=5.6cm,xmode=log,xmin=5,xmax=1000,ymin=0.08,ymax=0.22,
  xlabel={$Q/\mathrm{GeV}$},ylabel={$\alpha_s(Q)$},
  legend style={font=\scriptsize,draw=none,fill=none,at={(0.97,0.97)},anchor=north east},
  legend cell align=left,
  tick label style={font=\scriptsize},label style={font=\small},title={\small (a)},
  yticklabel style={/pgf/number format/fixed,/pgf/number format/precision=2}]
\addplot[domain=5:1000,samples=120,thick,dashed,blue!60!black]{0.1180/(1+(23/3)*0.1180/(4*pi)*ln(x^2/91.19^2))};
\addlegendentry{单圈}
\addplot[domain=5:1000,samples=120,thick,red!60!black]{4*pi/((23/3)*ln(x^2/0.2262^2))*(1-(116/3)/(23/3)^2*ln(ln(x^2/0.2262^2))/ln(x^2/0.2262^2))};
\addlegendentry{两圈（截断式）}
\addplot[only marks,mark=*,mark size=1.6pt,error bars/.cd,y dir=both,y explicit] coordinates {(91.19,0.1180) +- (0,0.0009)};
\node[font=\scriptsize,anchor=south west] at (axis cs:100,0.121) {$\alpha_s(M_Z)=0.1180$};
\node[font=\scriptsize,anchor=south west] at (axis cs:6,0.085) {$n_f=5$};
\end{axis}
\begin{axis}[at={(A.east)},anchor=west,xshift=1.45cm,width=0.47\textwidth,height=5.6cm,xmin=1.6,xmax=6.95,ymin=4.5,ymax=12,
  xlabel={$\ln(Q/\mathrm{GeV})$},ylabel={$1/\alpha_s(Q)$},tick label style={font=\scriptsize},label style={font=\small},title={\small (b)}]
\addplot[domain=1.609:6.908,samples=2,thick,blue!60!black]{1/0.1180+(23/3)/(2*pi)*(x-ln(91.19))};
% slope triangle attached to the line (values of the line at x=3.2 and x=4.2)
\draw[thin] (axis cs:3.2,6.8726) -- (axis cs:4.2,6.8726) -- (axis cs:4.2,8.0927);
\node[font=\scriptsize,anchor=north] at (axis cs:3.7,6.8726) {$1$};
\node[font=\scriptsize,anchor=west] at (axis cs:4.25,7.45) {斜率 $\beta_0/(2\pi)=1.22$};
\node[font=\scriptsize,align=left,anchor=north west] at (axis cs:1.75,11.7) {渐近自由：\\$Q\to\infty$ 时 $\alpha_s\to0$};
\end{axis}
\end{tikzpicture}
\caption{跑动耦合 $\alpha_s(Q)$（$n_f=5$，输入 PDG 2024 世界平均 $\alpha_s(M_Z)=0.1180\pm0.0009$\cite[式~(9.25)]{ParticleDataGroup:2024cfk}）。
(a) 单圈与两圈（截断式，$\Lambda^{(5)}_{\MSbar}=0.226\,\mathrm{GeV}$，在 $5\le Q\le1000\,\mathrm{GeV}$ 内与精确两圈跑动相差不到 $0.6\%$，VF-QCD-7）跑动；
(b) 单圈下 $1/\alpha_s$ 对 $\ln Q$ 为直线，斜率 $\beta_0/(2\pi)>0$：$Q$ 越大耦合越弱，这就是渐近自由。
图中全程取 $n_f=5$；$Q<m_b$ 与 $Q>m_t$ 处应分别改用 $n_f=4$ 与 $n_f=6$（阈值匹配），图中未画出。}
\label{fig:qcd-running-coupling}
\end{figure}
%%% END fig:qcd-running-coupling


%%% BEGIN fig:qcd-antiscreening
\begin{figure}[htbp]
\centering
\begin{tikzpicture}[fd,
  fa dip/.pic={\draw[line width=0.5pt,fill=gray!10] (0,0) ellipse[x radius=0.34,y radius=0.155];
               \node[font=\scriptsize,inner sep=0pt] at (-0.17,0) {$-$};
               \node[font=\scriptsize,inner sep=0pt] at (0.17,0) {$+$};},
  fa src/.style={circle,draw,line width=0.7pt,minimum size=0.62cm,inner sep=0pt,font=\small},
  fa lab/.style={font=\footnotesize},
  fa cap/.style={font=\footnotesize}]
  % ================= (a) QED =================
  \begin{scope}
    \foreach \a in {0,45,...,315} \pic[rotate=\a] at (\a:1.02) {fa dip};
    \node[fa src,fill=red!12] at (0,0) {$+$};
    \draw[lcaxis,line width=0.6pt] (22.5:0.42) -- (22.5:1.9) node[fa lab,anchor=south west,inner sep=1pt] {$r$ 增大};
    \node[fa lab,align=center] at (0,-1.85) {虚 $e^+e^-$ 对的“$-$”端靠近电荷：\\有效电荷随 $r$ 增大而减小（屏蔽）};
    \node[fa cap] at (0,-2.7) {(a)\ QED};
  \end{scope}
  % ================= (b) QCD =================
  \begin{scope}[xshift=5.3cm]
    \foreach \a in {0,90,180,270} \pic[rotate=\a] at (\a:1.02) {fa dip};
    \foreach \a/\c in {45/red!70!black,135/green!50!black,225/blue!70!black,315/red!70!black} {
      \draw[gluon,green!45!black] (\a:0.42) -- (\a:1.45);
      \fill[\c] (\a:1.56) circle[radius=0.075];
    }
    \node[fa src,fill=gray!15] at (0,0) {$\mathbf T$};
    \pic at (-1.6,-1.62) {fa dip};
    \node[fa lab,anchor=west] at (-1.15,-1.62) {夸克对：屏蔽，$-\tfrac43T_Fn_f$};
    \draw[gluon,green!45!black] (-1.92,-2.12) -- (-1.36,-2.12);
    \fill[red!70!black] (-1.27,-2.12) circle[radius=0.075];
    \node[fa lab,anchor=west] at (-1.15,-2.12) {胶子：色荷外移，反屏蔽，$+\tfrac{11}3C_A$};
    \node[fa cap] at (0,-2.7) {(b)\ QCD};
  \end{scope}
  % ================= (c) inset: effective coupling vs distance =================
  \begin{scope}[xshift=9.55cm,yshift=-1.3cm]
    \draw[lcaxis] (0,0) -- (3.6,0) node[fa lab,below left=1pt and -2pt] {$r$（对数刻度）};
    \draw[lcaxis] (0,0) -- (0,2.6) node[fa lab,anchor=south] {有效耦合};
    % schematic: QCD rises steeply towards large r, QED falls slowly; QCD stays above QED
    \draw[line width=0.9pt,blue!60!black] plot[smooth] coordinates {(0.15,0.78) (1.1,0.92) (2.0,1.2) (2.75,1.62) (3.3,2.2)};
    \draw[line width=0.9pt,red!60!black] plot[smooth] coordinates {(0.15,0.52) (1.2,0.44) (2.2,0.39) (3.3,0.35)};
    \node[fa lab,text=blue!60!black,anchor=south east] at (2.55,1.55) {QCD};
    \node[fa lab,text=red!60!black,anchor=south] at (2.3,0.42) {QED};
    \node[fa cap] at (1.8,-1.4) {(c)\ 有效耦合与距离（示意）};
  \end{scope}
\end{tikzpicture}
\caption{屏蔽与反屏蔽的直观图像（示意）。(a) QED：真空中的虚电子—正电子对屏蔽电荷，有效耦合在远处变小；
(b) QCD：夸克对同样屏蔽（$-\frac43T_Fn_f$），而带色的胶子把色荷“涂抹”到更大的区域，在短距离处看到的色荷变小（反屏蔽，$+\frac{11}3C_A$）；
(c) 两种理论的有效耦合随距离的变化趋势。二者之和 $\beta_0>0$（SU(3) 中 $n_f\le16$）。
这只是启发式图像：规范无关的推导见\cref{prop:prelim-beta}，把反屏蔽解释为胶子自旋的顺磁效应见\cite{Nielsen:1980sx}。}
\label{fig:qcd-antiscreening}
\end{figure}
%%% END fig:qcd-antiscreening


%%% BEGIN fig:qcd-ir-safety
\begin{figure}[htbp]
\centering
\begin{tikzpicture}[fd,
  fa lab/.style={font=\footnotesize},
  fa note/.style={font=\footnotesize,align=center},
  fa cap/.style={font=\footnotesize},
  fa hint/.style={draw=gray!60!black,line width=0.6pt,-{Stealth[length=4pt,width=3.4pt]}},
  fa ok/.style={draw=green!45!black,line width=0.8pt,rounded corners=2pt,fill=green!6,font=\footnotesize,align=center,inner sep=4pt},
  fa bad/.style={draw=red!60!black,line width=0.8pt,rounded corners=2pt,fill=red!5,font=\footnotesize,align=center,inner sep=4pt}]
  % ===================== (a) soft and collinear emission =====================
  \begin{scope}
    \node[hard blob,circle,minimum size=0.7cm,inner sep=0pt] (H) at (0,0) {$H$};
    \coordinate (V) at (1.75,0);
    \draw[quark] (H.east) -- (V);
    \draw[quark] (V) -- (3.75,0);
    \draw[gluon] (V) -- ++(36:1.75) coordinate (G);
    \node[vertex] at (V) {};
    \path[mom'={$p+k$}] (H.east) -- (V);
    \path[mom'={$p$},mom pos=0.62] (V) -- (3.75,0);
    \path[mom={$k$},mom sep=0.3cm,mom pos=0.42] (V) -- (G);
    \draw[line width=0.5pt] ($(V)+(0.55,0)$) arc[start angle=0,end angle=36,radius=0.55];
    \node[fa lab] at ($(V)+(17:0.75)$) {$\theta$};
    \draw[fa hint] ($(V)+(30:1.22)$) arc[start angle=30,end angle=6,radius=1.22];
    \node[fa lab,anchor=west] at ($(V)+(13:1.3)$) {共线：$\theta\to0$};
    \node[fa lab,anchor=south west,inner sep=1pt] at ($(G)+(0.04,0.02)$) {软：$E_k\to0$};
    \node[fa note] at (1.85,-1.15) {$\dfrac{1}{(p+k)^2}=\dfrac{1}{2E_pE_k(1-\cos\theta)}$};
    \node[fa note] at (1.85,-2.1) {发射概率 $\propto\dfrac{\dd E_k}{E_k}\,\dfrac{\dd\theta^2}{\theta^2}$};
    \node[fa cap] at (1.85,-2.95) {(a)\ 无质量夸克的软发射与共线发射};
  \end{scope}
  % ===================== (b) real-virtual cancellation =====================
  \begin{scope}[xshift=6.45cm,yshift=0.15cm]
    % real emission: the gluon crosses the cut
    \coordinate (L) at (0,0); \coordinate (R) at (2.2,0);
    \draw[photon] (-0.55,0) -- (L); \draw[photon] (R) -- (2.75,0);
    \draw[quark=0.13] (L) -- (0.42,0.55) -- (1.78,0.55) -- (R);
    \draw[quark=0.13] (R) -- (1.78,-0.55) -- (0.42,-0.55) -- (L);
    \draw[gluon] (0.6,0.55) .. controls (0.72,1.12) and (1.48,1.12) .. (1.6,0.55);
    \foreach \p in {(L),(R),(0.6,0.55),(1.6,0.55)} \node[vertex] at \p {};
    \draw[cut] (1.1,1.2) -- (1.1,-0.85);
    \node[fa note] at (1.1,-1.5) {实发射\\$N_\varepsilon\Big[\dfrac{2}{\varepsilon^2}+\dfrac{3}{\varepsilon}+\cdots\Big]$};
    \node[font=\large] at (3.3,0) {$+$};
    % virtual correction on one side of the cut, interfering with the Born amplitude
    \begin{scope}[xshift=4.4cm]
      \coordinate (L) at (0,0); \coordinate (R) at (2.2,0);
      \draw[photon] (-0.55,0) -- (L); \draw[photon] (R) -- (2.75,0);
      \draw[quark=0.13] (L) -- (0.42,0.55) -- (1.78,0.55) -- (R);
      \draw[quark=0.13] (R) -- (1.78,-0.55) -- (0.42,-0.55) -- (L);
      \draw[gluon] (0.66,0.55) -- (0.66,-0.55);
      \foreach \p in {(L),(R),(0.66,0.55),(0.66,-0.55)} \node[vertex] at \p {};
      \draw[cut] (1.1,1.2) -- (1.1,-0.85);
      \node[fa note] at (1.1,-1.5) {虚修正（及复共轭）\\$N_\varepsilon\Big[-\dfrac{2}{\varepsilon^2}-\dfrac{3}{\varepsilon}-\cdots\Big]$};
    \end{scope}
    \node[fa note] at (3.3,-2.45) {和 $=\dfrac32N_\varepsilon\ \xrightarrow{\ \varepsilon\to0\ }\ \dfrac34\,C_F\,\dfrac{\alpha_s}{\pi}$（有限）};
    \node[fa cap] at (3.3,-3.1) {(b)\ $e^+e^-$ 总截面中实发射与虚修正的相消};
  \end{scope}
  % ===================== (c) infrared safety =====================
  \begin{scope}[yshift=-5.25cm,
      fa prod/.style={circle,fill=gray!35,draw,line width=0.6pt,inner sep=0pt,minimum size=0.3cm},
      fa out/.style={draw,line width=0.8pt,-{Stealth[length=4.5pt,width=3.8pt]}}]
    % n-parton final state: two back-to-back partons from the production vertex
    \begin{scope}[xshift=0.9cm]
      \draw[fa out] (0,0) -- (1.05,0.36);  \draw[fa out] (0,0) -- (-1.05,-0.36);
      \node[fa prod] at (0,0) {};
      \node[fa lab] at (0,-0.8) {$n$ 个部分子};
    \end{scope}
    \draw[fa hint,line width=0.8pt] (2.25,0.18) -- (3.05,0.55);
    \draw[fa hint,line width=0.8pt] (2.25,-0.18) -- (3.05,-0.55);
    % + soft gluon emitted from one of the partons
    \begin{scope}[xshift=4.2cm,yshift=0.68cm]
      \draw[fa out] (0,0) -- (1.05,0.36);  \draw[fa out] (0,0) -- (-1.05,-0.36);
      \draw[gluon,green!45!black] (0.52,0.178) -- ++(108:0.78);
      \node[fa prod] at (0,0) {};
      \node[fa lab,text=green!35!black,anchor=west] at ($(0.52,0.178)+(108:0.62)+(0.12,0)$) {软胶子};
    \end{scope}
    % collinear splitting of one parton
    \begin{scope}[xshift=4.2cm,yshift=-0.72cm]
      \draw[plain] (0,0) -- (0.5,0.171);
      \draw[fa out] (0.5,0.171) -- (1.06,0.55);  \draw[fa out] (0.5,0.171) -- (1.13,0.26);
      \draw[fa out] (0,0) -- (-1.05,-0.36);
      \node[fa prod] at (0,0) {};
      \node[fa lab,anchor=west] at (1.2,0.38) {共线劈裂};
    \end{scope}
    \node[fa lab] at (4.2,-1.55) {$n+1$ 个部分子};
    % verdicts
    \node[fa ok,anchor=west] (ok) at (7.6,0.45) {红外安全：$O_{n+1}\to O_n$\\（如推力、Sterman--Weinberg 截面）};
    \node[fa bad,anchor=west] (bad) at (7.6,-0.75) {末态部分子数};
    \draw[red!70!black,line width=1.1pt] ($(bad.east)+(0.2,-0.17)$) -- ++(0.34,0.34) ($(bad.east)+(0.2,0.17)$) -- ++(0.34,-0.34);
    \node[fa lab,text=red!60!black,anchor=west] at ($(bad.east)+(0.62,0)$) {不安全};
    \node[fa cap] at (6.3,-2.15) {(c)\ 红外安全可观测量在软发射与共线劈裂下取值不变};
  \end{scope}
\end{tikzpicture}
\caption{红外发散与红外安全（示意）。(a) 无质量夸克发射胶子：传播子在软（$E_k\to0$）与共线（$\theta\to0$）极限下发散；
(b) 对足够单举的量，实发射与虚修正的发散相消（$e^+e^-$ 总截面，VF-EE-6；$N_\varepsilon$ 为两者共同的归一化因子，以 Born 截面 $\sigma_0$ 为单位，见\cref{prop:ee-oneloop}）；
(c) 红外安全可观测量在软发射与共线劈裂下取值不变，因而可以逐阶微扰计算（精确定义见\cref{def:ir-safe}）。}
\label{fig:qcd-ir-safety}
\end{figure}
%%% END fig:qcd-ir-safety


%%% BEGIN fig:qcd-gauge-orbit
\begin{figure}[htbp]
\centering
\begin{tikzpicture}[fd,
  fa orbit/.style={draw=blue!55!black,line width=0.7pt},
  fa gf/.style={draw=orange!80!black,line width=1.1pt},
  fa pt/.style={circle,fill=black,draw=none,inner sep=0pt,minimum size=4.2pt},
  fa lab/.style={font=\footnotesize}]
  % configuration space of gauge fields
  \fill[gray!8,rounded corners=8pt] (-0.45,-1.85) rectangle (10.75,5.05);
  \node[fa lab,anchor=north west,text=gray!45!black] at (-0.35,4.95) {规范场位形空间 $\{A^a_\mu(x)\}$};
  % gauge orbits {A^U}: x = x_i + a sin(0.9 y + phi)
  \draw[fa orbit] plot[domain=0:4.2,samples=40,variable=\t] ({1.8+0.18*sin(deg(0.9*\t))},{\t});
  \draw[fa orbit] plot[domain=0:4.2,samples=40,variable=\t] ({3.4+0.22*sin(deg(0.9*\t+0.6))},{\t});
  \draw[fa orbit] plot[domain=0:4.2,samples=40,variable=\t] ({5.0+0.22*sin(deg(0.9*\t+1.2))},{\t});
  \draw[fa orbit] plot[domain=0:4.2,samples=40,variable=\t] ({6.6+0.2*sin(deg(0.9*\t+1.8))},{\t});
  \draw[fa orbit] plot[domain=0:4.2,samples=40,variable=\t] ({8.4+0.18*sin(deg(0.9*\t+2.4))},{\t});
  \node[fa lab,text=blue!50!black,anchor=south] at (4.8,4.22) {规范轨道 $\{A^U\}$};
  % gauge-fixing surface F(A)=0 (folds back far from A=0)
  \draw[fa gf] (-0.2,1.02) -- (0.9,1.25) .. controls (3.3,1.75) and (5.5,1.45) .. (7.4,1.85)
               .. controls (9.3,2.25) and (9.65,3.45) .. (7.95,3.75);
  \node[fa lab,text=orange!65!black,anchor=south west,align=left] at (-0.3,1.32) {规范固定曲面\\$F(A)=0$};
  % intersections (computed from the curves above)
  \node[fa pt] at (1.973,1.427) {};
  \node[fa pt] at (3.601,1.558) {};
  \node[fa pt] at (5.105,1.615) {};
  \node[fa pt] at (6.561,1.716) {};
  \node[fa pt,red!65!black] at (8.235,2.116) {};
  \node[fa pt,red!65!black] at (8.300,3.667) {};
  \node[fa lab,anchor=north west,inner sep=1.5pt] at (2.0,1.38) {$A=0$};
  % Faddeev--Popov Jacobian: gauge direction along an orbit at the crossing
  \draw[lcaxis,line width=0.7pt,blue!55!black] (5.105,1.615) -- ++(-0.12,0.85);
  \node[fa lab,anchor=west,text=blue!50!black] at (5.02,2.42) {$\delta\theta$};
  % Gribov copies
  \draw[red!60!black,line width=0.6pt,dashed] (8.27,2.89) ellipse[x radius=0.36,y radius=1.12];
  \node[fa lab,text=red!55!black,anchor=west,align=left] at (9.2,2.9) {Gribov\\复本};
  % explanations in the lower band
  \node[fa lab,align=center,anchor=north] at (1.75,-0.12) {微扰论在 $A=0$ 附近展开：\\每条轨道只交一次};
  \node[fa lab,align=center,anchor=north] at (5.5,-0.12) {交点处的 Jacobian\\$\Delta_{\mathrm{FP}}=\det\dfrac{\delta F(A^U)}{\delta\theta}$\\给出鬼场};
  \node[fa lab,align=center,anchor=north,text=red!55!black] at (9.15,-0.12) {大场强区：同一轨道\\被交多次（Gribov 复本，\\微扰论中不出现）};
\end{tikzpicture}
\caption{规范固定与 Faddeev--Popov 行列式（示意）：规范固定曲面与每条规范轨道相交一次，交点处的 Jacobian 给出鬼场。
在远离 $A=0$ 的大场强区域，规范固定曲面可能与同一条轨道相交多次（Gribov 复本，红点）；它们在微扰论的任何一阶中都不出现。}
\label{fig:qcd-gauge-orbit}
\end{figure}
%%% END fig:qcd-gauge-orbit


%%% BEGIN fig:qcd-colour-flow
\begin{figure}[htbp]
\centering
\begin{tikzpicture}[fd,
  fa cl/.style={draw,line width=0.8pt,rounded corners=4pt,fd arrow=#1},
  fa cl/.default=0.5,
  fa idx/.style={font=\footnotesize,inner sep=1.5pt},
  fa eq/.style={font=\normalsize},
  fa cap/.style={font=\footnotesize}]
  % ===================== (a) gluon = colour--anticolour pair =====================
  \begin{scope}
    \draw[gluon] (0,0) -- (1.8,0);
    \node[fa idx,below] at (0.9,-0.18) {$a$};
    \node[fa eq] at (2.35,0) {$\leftrightarrow$};
    \draw[fa cl] (2.9,0.13) -- (4.7,0.13);
    \draw[fa cl] (4.7,-0.13) -- (2.9,-0.13);
    \node[fa cap,align=center] at (2.35,-1.05) {(a)\ 胶子 $\approx$ 一对方向相反的色线\\（双线记号；$-\tfrac1N$ 项见 (b)）};
  \end{scope}
  % ===================== (b) Fierz identity =====================
  \begin{scope}[xshift=6.5cm]
    % t^a_{ij} t^a_{kl}: upper quark line j -> i, lower quark line l -> k
    \draw[quark=0.25] (0,0.5) node[fa idx,left]{$j$} -- (1.6,0.5) node[fa idx,right]{$i$};
    \draw[quark=0.25] (1.6,-0.5) node[fa idx,right]{$l$} -- (0,-0.5) node[fa idx,left]{$k$};
    \draw[gluon] (0.8,0.5) -- (0.8,-0.5);
    \node[fa idx,right] at (0.95,0) {$a$};
    \foreach \y in {0.5,-0.5} \node[vertex] at (0.8,\y) {};
    \node[fa eq] at (2.35,0) {$=\ \dfrac12$};
    % rerouted colour lines: j -> k and l -> i
    \begin{scope}[xshift=3.3cm]
      \draw[fa cl] (0,0.5) node[fa idx,left]{$j$} -- (0.62,0.5) -- (0.62,-0.5) -- (0,-0.5) node[fa idx,left]{$k$};
      \draw[fa cl] (1.6,-0.5) node[fa idx,right]{$l$} -- (0.98,-0.5) -- (0.98,0.5) -- (1.6,0.5) node[fa idx,right]{$i$};
    \end{scope}
    \node[fa eq] at (5.75,0) {$-\ \dfrac{1}{2N}$};
    \begin{scope}[xshift=6.75cm]
      \draw[fa cl] (0,0.5) node[fa idx,left]{$j$} -- (1.6,0.5) node[fa idx,right]{$i$};
      \draw[fa cl] (1.6,-0.5) node[fa idx,right]{$l$} -- (0,-0.5) node[fa idx,left]{$k$};
    \end{scope}
    \node[font=\small] at (4.2,-1.1) {$t^a_{ij}\,t^a_{kl}=\tfrac12\big(\delta_{il}\delta_{kj}-\tfrac1N\,\delta_{ij}\delta_{kl}\big)$};
    \node[fa cap] at (4.2,-1.65) {(b)\ Fierz 恒等式};
  \end{scope}
  % ===================== (c) C_F from the Fierz identity, and C_A =====================
  \begin{scope}[yshift=-3.35cm,xshift=1.0cm]
    % quark self-energy colour factor
    \draw[quark=0.15] (0,0) -- (2.0,0);
    \draw[gluon] (0.5,0) .. controls (0.6,0.78) and (1.4,0.78) .. (1.5,0);
    \foreach \x in {0.5,1.5} \node[vertex] at (\x,0) {};
    \node[fa eq] at (2.6,0.05) {$=\ \dfrac12$};
    \begin{scope}[xshift=3.2cm]
      \draw[fa cl] (0,0) -- (2.0,0);
      \draw[fa cl=0.25] (1.0,0.45) circle[radius=0.28];
      \node[fa idx] at (1.0,0.45) {$N$};
    \end{scope}
    \node[fa eq] at (5.95,0.05) {$-\ \dfrac{1}{2N}$};
    \draw[fa cl] (6.75,0) -- (8.75,0);
    \node[fa eq] at (9.4,0.05) {$=\ C_F$};
    \draw[quark] (10.0,0) -- (12.0,0);
    \node[font=\small] at (6.0,-0.75) {$t^at^a=\tfrac12\big(N-\tfrac1N\big)\mathbb 1=C_F\,\mathbb 1$，\quad $C_F=\dfrac{N^2-1}{2N}$};
    % gluon self-energy colour factor
    \begin{scope}[yshift=-2.0cm,xshift=1.25cm]
      \draw[gluon] (0,0) -- (0.85,0);
      \draw[gluon] (0.85,0) arc[start angle=180,end angle=0,radius=0.5];
      \draw[gluon] (1.85,0) arc[start angle=0,end angle=-180,radius=0.5];
      \draw[gluon] (1.85,0) -- (2.7,0);
      \foreach \x in {0.85,1.85} \node[vertex] at (\x,0) {};
      \node[fa idx,below] at (0,-0.15) {$a$};
      \node[fa idx,below] at (2.7,-0.15) {$b$};
      \node[fa eq] at (3.4,0) {$=\ C_A$};
      \draw[gluon] (4.05,0) -- (6.0,0);
      \node[fa idx,below] at (4.05,-0.15) {$a$};
      \node[fa idx,below] at (6.0,-0.15) {$b$};
      \node[font=\small,anchor=west] at (6.5,0) {$f^{acd}f^{bcd}=C_A\,\delta^{ab}$，\quad $C_A=N$};
    \end{scope}
    \node[fa cap] at (6.0,-3.0) {(c)\ 色因子 $C_F$（夸克自能）与 $C_A$（胶子圈）};
  \end{scope}
\end{tikzpicture}
\caption{色流（双线）记号与 Fierz 恒等式；$C_F=\frac{N^2-1}{2N}$ 与 $C_A=N$ 的图示（VF-QCD-1）。单线箭头表示色指标的流向（夸克线即其费米子数流向），
闭合色线圈给出因子 $\delta_{ii}=N$；(c) 的第一行即用 (b) 把夸克自能的色因子 $t^at^a$ 化为 $C_F$。}
\label{fig:qcd-colour-flow}
\end{figure}
%%% END fig:qcd-colour-flow


%%% BEGIN fig:prelim-wilson-rules
\begin{figure}[htbp]
\centering
\begin{tikzpicture}[fd,
  fa end/.style={circle,draw,fill=white,line width=0.8pt,inner sep=0pt,minimum size=5.5pt},
  fa idx/.style={font=\footnotesize,inner sep=1.5pt},
  fa prop/.style={font=\small,inner sep=2pt},
  fa cap/.style={font=\footnotesize}]
  % ===================== (a) past-pointing line W[0,-infinity n] =====================
  \begin{scope}
    \draw[wilson=0.06] (0,0) -- (12.3,0);
    \node[fa idx,left] at (-0.05,0) {$-\infty n$};
    \node[fa end] at (12.3,0) {};
    \node[fa idx,right] at (12.4,0) {$0$};
    \foreach \x/\j in {1.5/1,5.1/2,8.7/3} {
      \draw[gluon] (\x,1.35) node[fa idx,above]{$\mu_\j,a_\j$} -- (\x,0);
      \path[mom={$k_\j$},mom sep=0.3cm] (\x,1.35) -- (\x,0.12);
      \node[vertex] at (\x,0) {};
    }
    \node[fa prop,below] at (0.75,-0.18) {$\ell=0$};
    \node[fa idx,anchor=south east] at (1.38,0.13) {$\ii g\,t^{a_1}n^{\mu_1}$};
    \node[fa prop,below] at (3.3,-0.1) {$\dfrac{\ii}{n\cdot k_1+\ii0}$};
    \node[fa prop,below] at (6.9,-0.1) {$\dfrac{\ii}{n\cdot(k_1+k_2)+\ii0}$};
    \node[fa prop,below] at (10.5,-0.1) {$\dfrac{\ii}{n\cdot(k_1+k_2+k_3)+\ii0}$};
    \node[fa cap] at (6.15,-1.55) {(a)\ 过去指向的 $W[0,-\infty n]$：$k_1$ 最靠近 $-\infty$，颜色因子 $t^{a_3}t^{a_2}t^{a_1}$};
  \end{scope}
  % ===================== (b) future-pointing line W[infinity n,0] =====================
  \begin{scope}[yshift=-4.1cm]
    \draw[wilson=0.94] (0,0) -- (12.3,0);
    \node[fa end] at (0,0) {};
    \node[fa idx,left] at (-0.1,0) {$0$};
    \node[fa idx,right] at (12.35,0) {$+\infty n$};
    \foreach \x/\j in {3.6/1,7.2/2,10.8/3} {
      \draw[gluon] (\x,1.35) node[fa idx,above]{$\mu_\j,a_\j$} -- (\x,0);
      \path[mom={$k_\j$},mom sep=0.3cm] (\x,1.35) -- (\x,0.12);
      \node[vertex] at (\x,0) {};
    }
    \node[fa idx,anchor=south east] at (3.48,0.13) {$\ii g\,t^{a_1}n^{\mu_1}$};
    \node[fa prop,below] at (1.8,-0.1) {$\dfrac{-\ii}{n\cdot(k_1+k_2+k_3)-\ii0}$};
    \node[fa prop,below] at (5.4,-0.1) {$\dfrac{-\ii}{n\cdot(k_2+k_3)-\ii0}$};
    \node[fa prop,below] at (9.0,-0.1) {$\dfrac{-\ii}{n\cdot k_3-\ii0}$};
    \node[fa prop,below] at (11.55,-0.18) {$\ell=0$};
    \node[fa cap] at (6.15,-1.55) {(b)\ 未来指向的 $W[\infty n,0]$：$k_1$ 最靠近 $0$，颜色因子 $t^{a_3}t^{a_2}t^{a_1}$};
  \end{scope}
\end{tikzpicture}
\caption{半无限 Wilson 线的动量空间规则（$k$ 流入 Wilson 线；颜色矩阵按路径排序，靠近终点者在左）：(a) 过去指向，(b) 未来指向；双线上的箭头为路径取向，空心圆为有限端点 $0$。
统一规则（VF-WL-1）：每个顶点 $\ii g\,t^{a_j}n^{\mu_j}$（图中只标出第一个）；有限端点与顶点之间、相邻顶点之间的每一段给一个程函传播子 $\ii/(n\cdot\ell+\ii0)$，
$\ell$ 为沿路径取向流动的动量，伸向无穷远的一段 $\ell=0$、不给传播子。于是 (a) 中第 $j$ 个顶点之后的一段 $\ell=k_1+\dots+k_j$；
(b) 中第 $j$ 个顶点之前的一段 $\ell=-(k_j+\dots+k_3)$，传播子即 $-\ii/(n\cdot(k_j+\dots+k_3)-\ii0)$。}
\label{fig:prelim-wilson-rules}
\end{figure}
%%% END fig:prelim-wilson-rules


%%% BEGIN fig:prelim-lightcone
\begin{figure}[htbp]
\centering
\begin{tikzpicture}[fd,
  fa lc/.style={draw=gray!60!black,line width=0.6pt,dash pattern=on 3pt off 2pt},
  fa hyp/.style={draw=black,line width=0.8pt,dash pattern=on 0.8pt off 1.6pt},
  fa grid/.style={draw=gray!60,line width=0.5pt,dash pattern=on 0.8pt off 1.6pt},
  fa pt/.style={circle,fill=black,draw=none,inner sep=0pt,minimum size=4.2pt},
  fa boost/.style={draw=red!60!black,line width=0.9pt,-{Stealth[length=5pt,width=4.5pt]}},
  fa leg/.style={draw=gray!50!black,line width=0.45pt},
  fa lab/.style={font=\footnotesize},
  fa wlab/.style={font=\footnotesize,fill=white,inner sep=1pt},
  fa cap/.style={font=\footnotesize}]
  % ===================== (a) the (z,t) plane =====================
  \begin{scope}
    % world-band of a hadron moving fast along +z (slope dt/dz = 1/v, v = 0.83)
    \draw[hadron blob,line width=0.5pt] (1.35,-0.45) -- (1.82,-0.45) -- (3.99,2.15) -- (3.52,2.15) -- cycle;
    \draw[fa leg] (1.9,-0.36) -- (2.35,-0.62);
    \node[fa lab,anchor=north west,align=left,inner sep=1pt] at (2.3,-0.55) {沿 $+z$ 快速运动\\的强子世界带};
    % axes and light-cone axes x^pm = (t pm z)/sqrt2
    \draw[lcaxis] (-2.6,0) -- (4.45,0) node[fa lab,below] {$z$};
    \draw[lcaxis] (0,-0.7) -- (0,3.0) node[fa lab,left] {$t$};
    \draw[fa lc,-{Stealth[length=4pt,width=3.4pt]}] (-0.6,-0.6) -- (2.75,2.75) node[fa lab,anchor=south west,inner sep=1pt] {$x^+$};
    \draw[fa lc,-{Stealth[length=4pt,width=3.4pt]}] (0.6,-0.6) -- (-2.75,2.75) node[fa lab,anchor=south east,inner sep=1pt] {$x^-$};
    % invariant hyperbola t^2 - z^2 = c and a boost of rapidity eta along it
    \draw[fa hyp] plot[domain=-1.2:1.2,samples=60,variable=\t] ({1.1*sinh(\t)},{1.1*cosh(\t)});
    \node[fa lab,anchor=south east,inner sep=1pt] at (-0.2,2.12) {$t^2-z^2=c$};
    \draw[fa leg] (-0.62,2.12) -- (-0.8,1.4);
    \draw[fa boost] plot[domain=0.06:0.74,samples=20,variable=\t] ({1.1*sinh(\t)},{1.1*cosh(\t)});
    \node[fa pt] at (0,1.1) {};
    \node[fa pt] at (0.977,1.471) {};
    \node[fa lab,text=red!55!black,anchor=north west] at (0.55,1.12) {推动 $\eta$};
    \node[fa cap] at (0.9,-1.85) {(a)\ $(z,t)$ 平面};
  \end{scope}
  % ===================== (b) the (ln v^+, ln v^-) plane =====================
  \begin{scope}[xshift=6.9cm]
    \begin{scope}
      \clip (0.05,0.05) rectangle (4.3,3.45);
      % lines of constant v^+ v^-  (anti-diagonals, slope -1)
      \foreach \c in {1.4,2.6,3.8,5.0,6.2} \draw[fa lc] (\c-3.6,3.6) -- (\c+0.2,-0.2);
      % lines of constant rapidity y = (1/2) ln(v^+/v^-)  (diagonals, slope +1)
      \foreach \d in {-2.4,-1.2,0,1.2,2.4} \draw[fa grid] (-0.2,-0.2+\d) -- (4.5,4.5+\d);
    \end{scope}
    \draw[lcaxis] (0,0) -- (4.55,0) node[fa lab,below left=1pt and -2pt] {$\ln v^+$};
    \draw[lcaxis] (0,0) -- (0,3.7) node[fa lab,left] {$\ln v^-$};
    % the boost: +eta in ln v^+, -eta in ln v^-
    \draw[fa leg] (1.5,2.3) -- (2.6,2.3) -- (2.6,1.2);
    \node[fa wlab,anchor=south] at (2.05,2.33) {$+\eta$};
    \node[fa wlab,anchor=west] at (2.64,1.75) {$-\eta$};
    \node[fa pt] (P) at (1.5,2.3) {};
    \node[fa pt] (Q) at (2.6,1.2) {};
    \draw[fa boost,shorten >=2pt,shorten <=2pt] (P) -- (Q);
    \node[fa wlab,text=red!55!black,anchor=north east] at (1.95,1.68) {推动 $\eta$};
    \node[fa wlab,anchor=south east] at (1.44,2.36) {$(\ln v^+,\ln v^-)$};
    \node[fa wlab,anchor=north] at (2.6,1.1) {$(\ln v^++\eta,\ \ln v^--\eta)$};
    % legend for the two line families
    \draw[fa lc] (0.1,-0.75) -- (0.75,-0.75);
    \node[fa lab,anchor=west] at (0.8,-0.75) {等 $v^+v^-$（反对角线）};
    \draw[fa grid] (0.1,-1.15) -- (0.75,-1.15);
    \node[fa lab,anchor=west] at (0.8,-1.15) {等快度 $y=\frac12\ln(v^+/v^-)$（对角线）};
    \node[fa cap] at (2.2,-1.85) {(b)\ 对数光锥平面};
  \end{scope}
\end{tikzpicture}
\caption{光锥坐标与 Lorentz 推动。(a) $(z,t)$ 平面中的光锥轴 $x^\pm=(t\pm z)/\sqrt2$ 与沿 $+z$ 快速运动的强子世界带；推动沿不变双曲线移动点。
(b) 对数光锥平面中，快度为 $\eta$ 的推动 $v^\pm\to \ee^{\pm\eta}v^\pm$ 是沿反对角线的平移（$\ln v^+$ 增加 $\eta$，$\ln v^-$ 减少 $\eta$），等 $v^+v^-$ 线为反对角线（虚线），等快度线为对角线（点线）。}
\label{fig:prelim-lightcone}
\end{figure}
%%% END fig:prelim-lightcone


%%% BEGIN fig:sudakov-hjjs
\begin{figure}[htbp]
\centering
\begin{tikzpicture}[fd,
  fa lab/.style={font=\small},
  fa note/.style={font=\footnotesize},
  fa cap/.style={font=\footnotesize},
  fa eq/.style={font=\normalsize}]
  % ===================== (a) leading region (reduced diagram) =====================
  \begin{scope}
    \node[hard blob,circle,minimum size=1.0cm,inner sep=0pt] (H) at (0,0) {$H$};
    \node[jet blob,ellipse,minimum width=2.15cm,minimum height=0.66cm,rotate=28] (J1) at (2.8,1.55) {};
    \node[jetB blob,ellipse,minimum width=2.15cm,minimum height=0.66cm,rotate=-28] (J2) at (2.8,-1.55) {};
    \node[soft blob,circle,minimum size=0.62cm,inner sep=0pt] (S) at (3.4,0) {};
    \draw[photon] (-1.75,0) -- (H.west);
    \path[mom={$q$},mom sep=0.28cm] (-1.75,0) -- (H.west);
    % quark: H -> J1 -> out;  antiquark: out -> J2 -> H  (arrows = fermion-number flow)
    \draw[quark] (H.60) -- (J1.180);
    \draw[quark] (J1.0) -- ++(28:1.25) coordinate (E1);
    \draw[quark] ($(J2.0)+(-28:1.25)$) coordinate (E2) -- (J2.0);
    \draw[quark] (J2.180) -- (H.-60);
    % momenta of the outgoing quark and antiquark (both point away from H)
    \path[mom={$p_1$},mom pos=0.55] (J1.0) -- (E1);
    \path[mom'={$p_2$},mom pos=0.55] (J2.0) -- (E2);
    % scalar-polarized collinear gluons J_i -> H
    \draw[longit] (J1.246) to[bend left=8] (H.14);  \longitmark{(J1.246)}{(H.14)}
    \draw[longit] (J2.114) to[bend right=8] (H.-14); \longitmark{(J2.114)}{(H.-14)}
    % soft gluons, attached to the jets only
    \draw[gluon,green!45!black] (S.118) -- (J1.290);
    \draw[gluon,green!45!black] (S.62) -- (J1.320);
    \draw[gluon,green!45!black] (S.-118) -- (J2.70);
    \draw[gluon,green!45!black] (S.-62) -- (J2.40);
    \node[fa lab] at (2.55,2.27) {$J_1$};
    \node[fa lab] at (2.55,-2.27) {$J_2$};
    \node[fa lab,anchor=west] at (S.east) {$S$};
    \node[fa note,anchor=west] at ($(E1)+(0.05,0)$) {沿 $n_A$};
    \node[fa note,anchor=west] at ($(E2)+(0.05,0)$) {沿 $n_B$};
    \node[fa note,text=blue!55!black,align=center,anchor=east] at (0.55,-1.3) {纵向极化\\共线胶子};
    \draw[blue!55!black,line width=0.4pt] (0.5,-1.05) -- (0.98,-0.5);
    \node[fa cap] at (1.75,-3.0) {(a)\ 领头区域（约化图）};
  \end{scope}
  % ===================== (b) factorized form F = H J1 J2 / S =====================
  \begin{scope}[yshift=-5.1cm,xshift=-2.45cm]
    \node[fa eq] at (-0.75,0) {$F\;=\;H\;\times$};
    % ---- J1: quark jet along n_A, Wilson line along n_B ----
    \begin{scope}[xshift=0.45cm]
      \coordinate (O) at (0,0);
      \node[jet blob,ellipse,minimum width=1.3cm,minimum height=0.5cm,rotate=28] (J) at (28:1.4) {};
      \draw[quark] (O) -- (J.180);
      \draw[quark] (J.0) -- ++(28:0.7);
      \draw[wilson=0.22] (-28:2.1) coordinate (W) -- (O);
      \draw[gluon] (J.275) -- ($(O)!0.47!(W)$);
      \node[vertex] at ($(O)!0.47!(W)$) {};
      \node[opvertex] at (O) {};
      \node[fa note,anchor=west] at (W) {$n_B$};
      \node[fa lab] at (1.25,1.45) {$J_1$};
    \end{scope}
    \node[fa eq] at (3.15,0) {$\times$};
    % ---- J2: antiquark jet along n_B, Wilson line along n_A ----
    \begin{scope}[xshift=3.75cm]
      \coordinate (O) at (0,0);
      \node[jetB blob,ellipse,minimum width=1.3cm,minimum height=0.5cm,rotate=-28] (J) at (-28:1.4) {};
      \draw[quark] (J.180) -- (O);
      \draw[quark] ($(J.0)+(-28:0.7)$) -- (J.0);
      \draw[wilson=0.8] (O) -- (28:2.1) coordinate (W);
      \draw[gluon] (J.85) -- ($(O)!0.47!(W)$);
      \node[vertex] at ($(O)!0.47!(W)$) {};
      \node[opvertex] at (O) {};
      \node[fa note,anchor=west] at (W) {$n_A$};
      \node[fa lab] at (1.25,-1.45) {$J_2$};
    \end{scope}
    \node[fa eq] at (6.55,0) {$\Big/$};
    % ---- S: two Wilson lines meeting at a cusp ----
    \begin{scope}[xshift=7.05cm]
      \coordinate (O) at (0,0);
      \draw[wilson=0.84] (O) -- (28:2.1) coordinate (WA);
      \draw[wilson=0.16] (-28:2.1) coordinate (WB) -- (O);
      \node[soft blob,circle,minimum size=0.52cm,inner sep=0pt] (Sb) at (1.5,0) {};
      \draw[gluon,green!45!black] (Sb.110) -- ($(O)!0.5!(WA)$);
      \draw[gluon,green!45!black] (Sb.-110) -- ($(O)!0.5!(WB)$);
      \node[vertex] at ($(O)!0.5!(WA)$) {};
      \node[vertex] at ($(O)!0.5!(WB)$) {};
      \node[opvertex] at (O) {};
      \node[fa note,anchor=west] at (WA) {$n_A$};
      \node[fa note,anchor=west] at (WB) {$n_B$};
      \node[fa lab] at (1.25,1.45) {$S$};
    \end{scope}
    \node[fa cap] at (4.2,-1.85) {(b)\ 因子化形式 $F=H\,J_1J_2/S$};
  \end{scope}
\end{tikzpicture}
\caption{Sudakov 形状因子的领头区域与因子化 $F=H\,J_1J_2/S$（$\gamma^*\to q(p_1)\,\bar q(p_2)$，$p_1$ 沿 $n_A$、$p_2$ 沿 $n_B$）。
(a) 硬子图 $H$、两个喷注子图 $J_1$、$J_2$ 与只连接喷注的软子图 $S$（绿色）；Feynman 规范下 $J_i$ 与 $H$ 之间还有任意多条标量（纵向）极化的共线胶子（蓝色，以圆点标记）。
(b) 喷注函数中的 Wilson 线沿另一喷注方向，软函数是两条 Wilson 线的真空期望（$\otimes$ 为算符所在点，双线箭头与被替代的夸克线的费米子数流向一致）。}
\label{fig:sudakov-hjjs}
\end{figure}
%%% END fig:sudakov-hjjs


%%% BEGIN fig:sudakov-cusp
\begin{figure}[htbp]
\centering
\begin{tikzpicture}[fd,
  fa lc/.style={draw=gray!60!black,line width=0.55pt,dash pattern=on 3pt off 2pt},
  fa ext/.style={draw=gray!55!black,line width=0.5pt,dash pattern=on 1.5pt off 1.5pt},
  fa lab/.style={font=\small},
  fa note/.style={font=\footnotesize},
  fa cap/.style={font=\footnotesize},
  fa eq/.style={font=\normalsize}]
  % ===================== (a) the cusp: incoming along v2 from the past, outgoing along v1 =====================
  \begin{scope}
    \draw[lcaxis,gray!60!black] (-2.15,0) -- (2.3,0) node[fa note,below] {$z$};
    \draw[lcaxis,gray!60!black] (0,-2.05) -- (0,2.3) node[fa note,left] {$t$};
    \draw[fa lc] (-1.85,-1.85) -- (1.85,1.85);
    \draw[fa lc] (1.85,-1.85) -- (-1.85,1.85) node[fa note,anchor=south east,text=gray!45!black,inner sep=1pt] {光锥};
    % continuation of the incoming direction v2 beyond the cusp
    \draw[fa ext] (0,0) -- (68:1.45);
    \draw[wilson=0.5] (248:1.95) node[fa lab,anchor=north,inner sep=2pt] {$v_2$} -- (0,0);
    \draw[wilson=0.62] (0,0) -- (112:1.95) node[fa lab,anchor=south,inner sep=2pt] {$v_1$};
    \draw[line width=0.5pt] (68:0.85) arc[start angle=68,end angle=112,radius=0.85];
    \node[fa lab] at (90:1.08) {$\beta$};
    \node[opvertex] at (0,0) {};
    \node[fa cap,align=center] at (0.1,-2.62) {(a)\ 类时 Wilson 线的尖点，$\cosh\beta=v_1\cdot v_2$};
  \end{scope}
  % ===================== (b) light-like limit =====================
  \begin{scope}[xshift=5.6cm]
    \draw[fa lc] (0,0) -- (45:2.45);
    \draw[fa lc] (0,0) -- (-45:2.45);
    % in the limit both lines lie on the light cone (drawn on top of the dashed cone)
    \draw[wilson=0.5] (225:1.95) node[fa lab,anchor=north east,inner sep=1pt] {$v_2\to n_A$} -- (0,0);
    \draw[wilson=0.62] (0,0) -- (135:1.95) node[fa lab,anchor=south east,inner sep=1pt] {$v_1\to n_B$};
    \node[opvertex] at (0,0) {};
    \node[fa note,align=left,anchor=west] at (1.45,0.05)
      {$\beta\to\infty$：$\Gamma(\beta)\simeq\dfrac{\alpha_sC_F}{\pi}\,(\beta-1)$\\[3pt]
       线性依赖于 $\beta$，斜率即单圈\\[1pt]
       $\Gamma_{\mathrm{cusp}}=\dfrac{\alpha_sC_F}{\pi}=\dfrac{\alpha_s}{4\pi}\,\Gamma_0$};
    \node[fa cap] at (2.2,-2.62) {(b)\ 类光极限};
  \end{scope}
  % ===================== (c) one-loop graphs =====================
  \begin{scope}[yshift=-5.55cm,xshift=-0.6cm]
    % exchange between the two lines
    \begin{scope}
      \draw[wilson=0.17] (248:1.7) -- (0,0);
      \draw[wilson=0.83] (0,0) -- (112:1.7);
      \draw[gluon] (248:0.95) to[bend left=62] (112:0.95);
      \foreach \a in {248,112} \node[vertex] at (\a:0.95) {};
      \node[opvertex] at (0,0) {};
      \node[fa note] at (-0.35,-1.95) {交换};
    \end{scope}
    \node[fa eq] at (1.25,0) {$+$};
    % self-energy of line 1
    \begin{scope}[xshift=2.55cm]
      \draw[wilson=0.5] (248:1.7) -- (0,0);
      \draw[wilson=0.6] (0,0) -- (112:1.7);
      \draw[gluon] (112:0.42) .. controls ($(112:0.42)+(22:0.85)$) and ($(112:1.42)+(22:0.85)$) .. (112:1.42);
      \foreach \r in {0.42,1.42} \node[vertex] at (112:\r) {};
      \node[opvertex] at (0,0) {};
      \node[fa note] at (0,-1.95) {线 1 的自能};
    \end{scope}
    \node[fa eq] at (4.1,0) {$+$};
    % self-energy of line 2
    \begin{scope}[xshift=5.35cm]
      \draw[wilson=0.4] (248:1.7) -- (0,0);
      \draw[wilson=0.5] (0,0) -- (112:1.7);
      \draw[gluon] (248:1.42) .. controls ($(248:1.42)+(-22:0.85)$) and ($(248:0.42)+(-22:0.85)$) .. (248:0.42);
      \foreach \r in {0.42,1.42} \node[vertex] at (248:\r) {};
      \node[opvertex] at (0,0) {};
      \node[fa note] at (0,-1.95) {线 2 的自能};
    \end{scope}
    \node[fa eq,anchor=west] at (7.0,0) {$\Longrightarrow\ \ \Gamma(\beta)=\dfrac{\alpha_sC_F}{\pi}\,\big(\beta\coth\beta-1\big)$};
    \node[fa cap] at (6.0,-2.55) {(c)\ 单圈图：交换与两条线的自能（自能给出 $-1$，使 $\Gamma(0)=0$）};
  \end{scope}
\end{tikzpicture}
\caption{Wilson 线尖点及其单圈修正。(a) 类时 Wilson 线 $W[\infty v_1,0]\,W[0,-\infty v_2]$（$v_i^2=1$，$v_i^0>0$；双线箭头为路径方向）：沿 $v_2$ 自过去入射，在原点折向 $v_1$；
尖点角 $\beta$ 由 $\cosh\beta=v_1\cdot v_2$ 定义，$\beta=0$ 即一条不折的直线（图中所画角度为示意）。(b) 类光极限：$v_2$、$v_1$ 分别趋于 $n_A$、$n_B$ 方向，$\beta\to\infty$。
(c) 包含两线自能的单圈尖点反常维度为 $\Gamma(\beta)=\frac{\alpha_sC_F}\pi(\beta\coth\beta-1)$，类光极限下线性依赖于 $\beta$，斜率给出单圈 $\Gamma_{\rm cusp}$（VF-CONV-23）。}
\label{fig:sudakov-cusp}
\end{figure}
%%% END fig:sudakov-cusp


%%% BEGIN fig:sudakov-dipole
\begin{figure}[htbp]
\centering
\begin{tikzpicture}[fd,
  fa soft/.style={gluon,green!45!black},
  fa web/.style={soft blob,line width=0.7pt,pattern=north east lines,pattern color=green!45!black},
  fa over/.style={preaction={draw=white,line width=4.2pt,decorate=false}},
  fa lab/.style={font=\small},
  fa note/.style={font=\footnotesize},
  fa cap/.style={font=\footnotesize},
  fa legs/.pic={\foreach \a/\n in {66/1,22/2,-22/3,-66/4} {
                  \draw[plain] (\a:0.4) -- (\a:2.2) node[font=\small,anchor=\a+180,inner sep=2pt] {$\n$};}
                \node[hard blob,circle,minimum size=0.8cm,inner sep=0pt] at (0,0) {$H$};}]
  % ===================== (a) dipole terms =====================
  \begin{scope}
    \pic at (0,0) {fa legs};
    \draw[fa soft] (66:1.75) to[bend left=38] (22:1.75);
    \draw[fa soft] (-22:1.75) to[bend left=38] (-66:1.75);
    \foreach \a in {66,22,-22,-66} \node[vertex] at (\a:1.75) {};
    \node[fa note,anchor=south west,inner sep=1pt] at (44:2.12) {$\mathbf T_1\cdot\mathbf T_2$};
    \node[fa note,anchor=north west,inner sep=1pt] at (-44:2.12) {$\mathbf T_3\cdot\mathbf T_4$};
    \node[fa note,align=left,anchor=west] at (2.75,0) {对所有\\$i\neq j$ 求和};
    \node[fa cap] at (1.3,-3.05) {(a)\ 两圈以内：偶极子（色关联 $\mathbf T_i\cdot\mathbf T_j$）};
  \end{scope}
  % ===================== (b) three-loop webs (non-dipole) =====================
  % web blob in the free sector between legs 2 and 3; gluons that pass over a leg get a white gap
  \begin{scope}[xshift=7.2cm]
    \pic at (0,0) {fa legs};
    \draw[fa soft,fa over] (1.66,0.27) -- (66:1.5);
    \draw[fa soft] (1.86,0.2) -- (22:2.02);
    \draw[fa soft] (1.86,-0.2) -- (-22:2.02);
    \node[fa web,circle,minimum size=0.62cm,inner sep=0pt] (W) at (1.75,0) {};
    \foreach \p in {(66:1.5),(22:2.02),(-22:2.02)} \node[vertex] at \p {};
    \node[fa note] at (0.9,-2.6) {连接三条外线};
  \end{scope}
  \begin{scope}[xshift=11.6cm]
    \pic at (0,0) {fa legs};
    \draw[fa soft,fa over] (1.66,0.27) -- (66:1.5);
    \draw[fa soft,fa over] (1.66,-0.27) -- (-66:1.5);
    \draw[fa soft] (1.86,0.2) -- (22:2.02);
    \draw[fa soft] (1.86,-0.2) -- (-22:2.02);
    \node[fa web,circle,minimum size=0.62cm,inner sep=0pt] (W) at (1.75,0) {};
    \foreach \p in {(66:1.5),(22:2.02),(-22:2.02),(-66:1.5)} \node[vertex] at \p {};
    \node[fa note] at (0.9,-2.6) {连接四条外线};
  \end{scope}
  \node[fa cap] at (9.45,-3.05) {(b)\ 三圈起：非偶极子 web（阴影）};
  % ===================== anomalous dimension (VF-SUD-5) =====================
  \node[font=\small] at (5.6,-3.9)
    {$\boldsymbol\Gamma=\displaystyle\sum_{(i,j)}\frac{\mathbf T_i\cdot\mathbf T_j}{2}\,\gamma_{\rm cusp}(\alpha_s)\ln\frac{\mu^2}{-s_{ij}}+\sum_i\gamma^i(\alpha_s)$
     \quad（两圈以内成立）};
\end{tikzpicture}
\caption{多腿振幅的红外结构：(a) 两圈以内由偶极子项（色关联 $\mathbf T_i\cdot\mathbf T_j$）给出，反常维度如图下公式（VF-SUD-5，$(i,j)$ 取有序对）；
(b) 三圈起出现连接三条或四条外线的非偶极子 web（偶极子公式作为全阶公式在三圈失效，[V]）。}
\label{fig:sudakov-dipole}
\end{figure}
%%% END fig:sudakov-dipole


%%% BEGIN fig:ee-pi-plane
\begin{figure}[htbp]
\centering
\begin{tikzpicture}[fd,
  fa lab/.style={font=\footnotesize},
  fa contour/.style={draw=blue!65!black,line width=1.1pt,postaction={decorate,decoration={markings,
      mark=at position 0.13 with {\arrow{Stealth[length=6pt,width=5pt]}},
      mark=at position 0.5 with {\arrow{Stealth[length=6pt,width=5pt]}},
      mark=at position 0.87 with {\arrow{Stealth[length=6pt,width=5pt]}}}}},
  fa cutline/.style={draw=black,line width=1.0pt,decorate,decoration={zigzag,amplitude=1.7pt,segment length=4.6pt}},
  fa sm/.style={circle,draw=red!65!black,fill=white,line width=0.8pt,inner sep=0pt,minimum size=5.5pt}]
  \draw[lcaxis] (-4.0,0) -- (4.7,0) node[fa lab,anchor=west] {$\Real q^2$};
  \draw[lcaxis] (0,-3.15) -- (0,3.3) node[fa lab,left] {$\Imag q^2$};
  \draw[fa cutline] (0,0) -- (4.35,0);
  % contour of the contour formula: |q^2| = s, counter-clockwise from s+i0 to s-i0
  \draw[fa contour] (2.5:2.3) arc[start angle=2.5,end angle=357.5,radius=2.3];
  \node[fa lab,anchor=south west,inner sep=1.5pt] at (2.32,0.08) {$s+\ii0$};
  \node[fa lab,anchor=north west,inner sep=1.5pt] at (2.32,-0.08) {$s-\ii0$};
  \node[fa lab,text=blue!55!black,anchor=south east] at (135:2.38) {$|q^2|=s$};
  % smearing points s pm i Delta
  \draw[red!65!black,line width=0.5pt,dash pattern=on 1.5pt off 1.5pt] (2.3,-1.2) -- (2.3,1.2);
  \node[fa sm] at (2.3,1.2) {};
  \node[fa sm] at (2.3,-1.2) {};
  \node[fa lab,text=red!55!black,anchor=west] at (2.42,1.2) {$s+\ii\Delta$};
  \node[fa lab,text=red!55!black,anchor=west] at (2.42,-1.2) {$s-\ii\Delta$};
  % Euclidean point
  \fill (-1.3,0) circle[radius=2.2pt];
  \node[fa lab,anchor=north] at (-1.3,-0.08) {$-Q^2$};
  \node[fa lab,align=center,anchor=south] at (-1.3,0.18) {PT 可靠\\（欧氏区）};
  % cut and resonance region
  \draw[gray!60!black,line width=0.45pt] (3.9,0.1) -- (3.9,0.6);
  \node[fa lab,anchor=south,inner sep=1.5pt] at (3.9,0.6) {割线 $[0,\infty)$};
  \draw[gray!60!black,line width=0.5pt,-{Stealth[length=4pt,width=3.4pt]}] (0.82,-0.6) -- (0.36,-0.1);
  \node[fa lab,anchor=north west,align=left,inner sep=1pt] at (0.72,-0.56) {共振区\\（小 $s$）};
\end{tikzpicture}
\caption{流关联函数 $\Pi(q^2)$ 的解析结构：割线 $[0,\infty)$、欧氏点 $-Q^2$、围道公式 $R(s)=\frac1{2\pi}\int_{-\pi}^{\pi}\dd\theta\,D(s\ee^{\ii\theta})$ 的积分路径 $|q^2|=s$
（$D$ 的宗量 $Q^2=-q^2$ 走遍该圆周，在 $q^2$ 平面中即自 $s+\ii0$ 逆时针到 $s-\ii0$）以及抹平所用的点 $s\pm\ii\Delta$。在复 $q^2$ 处使用 OPE 是额外的假设（\cref{rem:ee-complex-ope}）。}
\label{fig:ee-pi-plane}
\end{figure}
%%% END fig:ee-pi-plane
